% ============================================================================
%  TEMPLATE LaTeX — Rodapé Institucional UFMT
%  Disciplina : Algoritmos 1
%  Professora : Rafaela Souza Francisco
%  Curso      : Bacharelado em Ciência da Computação — UFMT, Cuiabá
% ============================================================================
\documentclass[12pt, a4paper]{article}

% --- Codificação e idioma ---------------------------------------------------
\usepackage[utf8]{inputenc}
\usepackage[T1]{fontenc}
% \usepackage[brazil]{babel}  % descomente se tiver texlive-lang-portuguese

% --- Fonte ------------------------------------------------------------------
% \usepackage{lmodern}        % descomente se disponível no seu sistema

% --- Layout de página -------------------------------------------------------
\usepackage[
  top=2.5cm,
  bottom=3.0cm,       % espaço extra para o rodapé
  left=2.5cm,
  right=2.5cm
]{geometry}

% --- Pacotes necessários ----------------------------------------------------
\usepackage{graphicx}          % inserir imagens
\usepackage[brazil]{babel}
\usepackage{fancyhdr}          % cabeçalho/rodapé customizado
\usepackage{xcolor}            % cores
\usepackage{hyperref}          % links clicáveis
\usepackage{tikz}              % desenhos (usado na linha decorativa)
\usepackage{amsmath, amssymb}  % símbolos matemáticos (útil para Algoritmos)
\usepackage{lastpage}          % referência à última página

\usepackage{listings}
\usepackage{xcolor}
\renewcommand{\lstlistingname}{Código}

\lstset{
    basicstyle=\ttfamily\small,
    keywordstyle=\color{blue}\bfseries,
    commentstyle=\color{gray}\itshape,
    stringstyle=\color{red},
    numbers=left,
    numberstyle=\tiny\color{gray},
    numbersep=8pt,
    frame=single,
    breaklines=true,
    tabsize=4,
    showstringspaces=false
}

\usetikzlibrary{shapes.geometric, arrows}
\tikzstyle{iniciofim} = [ellipse, draw, minimum width=3cm, minimum height=1cm, text centered]
\tikzstyle{processo} = [rectangle, draw, minimum width=5cm, minimum height=1cm, text centered]
\tikzstyle{seta} = [->, thick]
\tikzstyle{decisao} = [diamond, minimum width=3.5cm, minimum height=2cm, text centered, draw=black]

% --- Cores institucionais UFMT ---------------------------------------------
\definecolor{ufmtAzul}{RGB}{0, 51, 102}       % azul escuro institucional
\definecolor{ufmtVerde}{RGB}{0, 105, 62}       % verde institucional
\definecolor{ufmtCinza}{RGB}{100, 100, 100}    % cinza para texto secundário

% --- Dados do documento (EDITE AQUI) ----------------------------------------
\newcommand{\NomeAluno}{Thalles Eduardo Silva Mota}             % ← Seu nome
\newcommand{\Disciplina}{Algoritmos 1}
\newcommand{\Professora}{Profª. Rafaela Souza Francisco}
\newcommand{\Curso}{Bacharelado em Ciência da Computação}
\newcommand{\Instituicao}{UFMT — Universidade Federal de Mato Grosso}
\newcommand{\Semestre}{2026/1}
\renewcommand{\abstractname}{Resumo}
% ← Semestre atual
% Se tiver o logo em arquivo, defina o caminho abaixo:
% \newcommand{\LogoPath}{logo_ufmt.png}

% ============================================================================
%  CONFIGURAÇÃO DO RODAPÉ  (o coração do template)
% ============================================================================
\pagestyle{fancy}
\fancyhf{}                     % limpa cabeçalho e rodapé padrão

% --- Cabeçalho (opcional — descomente se quiser) ----------------------------
% \fancyhead[L]{\small\color{ufmtCinza}\Disciplina}
% \fancyhead[R]{\small\color{ufmtCinza}\NomeAluno}
% \renewcommand{\headrulewidth}{0.4pt}

% --- Rodapé -----------------------------------------------------------------
\renewcommand{\footrulewidth}{0pt}  % remove linha padrão (usaremos tikz)

\fancyfoot[C]{%
  \vspace{2pt}%
  % ── Linha decorativa com gradiente ──────────────────────────────────
  \begin{tikzpicture}[remember picture, overlay]
    \draw[
      line width=1.2pt,
      ufmtAzul
    ] ([yshift=8pt]current page.south west |- 0,0) ++(2.5cm,0)
      -- ++(\dimexpr\textwidth\relax, 0);
    % Pequeno detalhe em verde (identidade visual)
    \fill[ufmtVerde]
      ([yshift=6.5pt]current page.south west |- 0,0) ++(2.5cm,0)
      rectangle ++(4cm, 1.8pt);
  \end{tikzpicture}%
  %
  % ── Conteúdo do rodapé ──────────────────────────────────────────────
  \begin{minipage}[t]{0.12\textwidth}
    \vspace{-2pt}
    % =====================================================================
    % LOGO: Descomente a linha abaixo e comente o \fbox quando tiver o logo
    % \includegraphics[height=1.1cm]{\LogoPath}
    % =====================================================================
    % Placeholder visual enquanto não tiver o .png:
    \includegraphics[height=1.1cm]{images/UFMT - logo texto.png}%
  \end{minipage}%
  \hfill
  \begin{minipage}[t]{0.60\textwidth}
    \centering
    \vspace{-2pt}
    {\footnotesize\color{ufmtAzul}\textbf{\Instituicao}}\\[1pt]
    {\scriptsize\color{ufmtCinza}%
      \Curso\ \textbullet\ \Disciplina\ — \Semestre}\\[1pt]
    {\scriptsize\color{ufmtCinza}\Professora}
  \end{minipage}%
  \hfill
  \begin{minipage}[t]{0.15\textwidth}
    \raggedleft
    \vspace{-2pt}
    {\footnotesize\color{ufmtAzul}\textbf{\thepage}\,/\,\pageref{LastPage}}
  \end{minipage}%
}

% Estilo da primeira página (pode ser diferente)
\fancypagestyle{plain}{%
  \fancyhf{}%
  \renewcommand{\headrulewidth}{0pt}%
  \fancyfoot[C]{%
    \vspace{2pt}%
    \begin{tikzpicture}[remember picture, overlay]
      \draw[line width=1.2pt, ufmtAzul]
        ([yshift=8pt]current page.south west |- 0,0) ++(2.5cm,0)
        -- ++(\dimexpr\textwidth\relax, 0);
      \fill[ufmtVerde]
        ([yshift=6.5pt]current page.south west |- 0,0) ++(2.5cm,0)
        rectangle ++(4cm, 1.8pt);
    \end{tikzpicture}%
    \begin{minipage}[t]{0.12\textwidth}
      \vspace{-2pt}
      \fbox{\footnotesize\textbf{\color{ufmtAzul}UFMT}}%
    \end{minipage}%
    \hfill
    \begin{minipage}[t]{0.60\textwidth}
      \centering\vspace{-2pt}
      {\footnotesize\color{ufmtAzul}\textbf{\Instituicao}}\\[1pt]
      {\scriptsize\color{ufmtCinza}%
        \Curso\ \textbullet\ \Disciplina\ — \Semestre}\\[1pt]
      {\scriptsize\color{ufmtCinza}\Professora}
    \end{minipage}%
    \hfill
    \begin{minipage}[t]{0.15\textwidth}
      \raggedleft\vspace{-2pt}
      {\footnotesize\color{ufmtAzul}\textbf{\thepage}\,/\,\pageref{LastPage}}
    \end{minipage}%
  }%
}

% ============================================================================
%  DOCUMENTO DE EXEMPLO
% ============================================================================
\begin{document}

\begin{center}
  {\Large\bfseries\color{ufmtAzul} Exercícios: Variáveis em C}\\[6pt]
  {\large \NomeAluno}\\[3pt]
  {\normalsize\color{ufmtCinza} \today}
\end{center}

\vspace{0.5cm}

\begin{abstract}
    \textbf{Este trabalho apresenta duas soluções para o problema clássico de 
    troca de valores entre duas variáveis inteiras na linguagem C. A primeira 
    abordagem utiliza uma variável auxiliar temporária para preservar um dos 
    valores durante a permutação. A segunda abordagem realiza a troca 
    exclusivamente por meio de operações aritméticas de soma e subtração, 
    dispensando memória adicional. Ambas as implementações demonstram conceitos 
    fundamentais de manipulação de variáveis, atribuição e expressões aritméticas 
    em programação estruturada.}
\end{abstract}

\noindent\textbf{Palavras-chave: troca de variáveis; variável auxiliar; 
operações aritméticas; linguagem C; atribuição} .

\section*{Introdução}
  A manipulação de variáveis constitui uma das operações mais elementares no 
  estudo de programação de computadores. Compreender como os dados são 
  armazenados, acessados e modificados na memória é um passo essencial para o 
  desenvolvimento do raciocínio lógico-computacional exigido na formação em 
  Ciência da Computação.


Dentre os problemas introdutórios frequentemente propostos, a troca de valores 
entre duas variáveis ocupa posição de destaque por sua simplicidade aparente e 
pela riqueza conceitual que oferece. Embora o enunciado seja direto — permutar 
o conteúdo(trocar os valores) de duas variáveis inteiras — sua resolução mobiliza 
conceitos fundamentais como declaração de variáveis, comando de atribuição, 
uso de variáveis auxiliares e construção de expressões aritméticas.


O presente trabalho aborda duas estratégias distintas para a solução desse 
problema na linguagem C. A primeira emprega uma variável temporária como espaço 
intermediário de armazenamento, garantindo que nenhum valor seja perdido durante 
o processo de troca. A segunda estratégia elimina a necessidade de memória 
adicional ao explorar propriedades das operações de soma e subtração para 
rearranjar os valores utilizando apenas as duas variáveis originais.


O objetivo é não apenas apresentar as implementações, mas também discutir as 
vantagens, limitações e os conceitos de programação envolvidos em cada abordagem, 
contribuindo para a consolidação dos fundamentos de algoritmos em uma programação 
computacional estruturada.
\pagebreak

\section*{Estratégia 1 - Troca com variável auxiliar}


A primeira solução utiliza uma variável auxiliar temporária chamada $aux$ para 
intermediar a troca. O processo ocorre em três passos sequenciais: primeiro, o 
valor de $a$ é copiado para temp; em seguida, a recebe o valor de $b$; por fim, 
$b$ recebe o valor armazenado em $aux$, que corresponde ao valor original de $a$.


Essa abordagem é considerada o método padrão de troca de variáveis por sua 
simplicidade e confiabilidade. O custo adicional de memória é mínimo — apenas 
uma variável inteira extra — e não há risco de erros numéricos, independentemente 
dos valores envolvidos. Tais conceitos, são demonstrados na seção 2.1.3 do livro 
citado:\cite{medina2006}.


O Código 1 apresenta a implementação completa desta solução.


\begin{lstlisting}[language=C, caption={Troca com variável auxiliar}, label={lst:q1}]
/* =====================================================
   Estrategia:
   - Usa uma terceira variavel (aux) para guardar
     temporariamente o valor de 'a' antes de sobrescreve-lo.
   - Passo 1: aux recebe o valor de a
   - Passo 2: a recebe o valor de b (a antigo esa salvo)
   - Passo 3: b recebe o valor de aux (que e o a antigo)
   ==================================================== */

#include <stdio.h>

int main() {
    int a, b, aux;

    printf("Digite dois inteiros: ");
    scanf("%d %d", &a, &b);

    printf("Antes: a = %d, b = %d\n", a, b);

    aux = a;
    a = b;
    b = aux;

    printf("Depois: a = %d, b = %d\n", a, b);

    return 0;
}
\end{lstlisting}



    

\section*{Estratégia 2 - Troca sem variável auxiliar}

A segunda solução elimina a necessidade de uma terceira variável ao explorar uma 
propriedade das operações de soma e subtração: se $a$ recebe $a+b$, então $a-b$ 
recupera o valor original de $a$, e uma segunda subtração recupera o valor 
original de $b$.


Concretamente, o processo se dá em três passos: $a = a+b$ faz com que $a$ passe 
a conter a soma dos dois valores; $b = a-b$ atribui a $b$. Ao final, os valores 
estão trocados sem uso de memória adicional.


Embora elegante, essa técnica apresenta uma limitação prática: se a soma de $a$ 
e $b$ ultrapassar o limite do tipo $int$(aproximadamente $\pm2,15 \times 10^9$ 
em em arquiteturas de 32 bits), ocorrerá *overflow*, produzindo resultados 
incorretos. Portanto, sua solução está proposta no seguinte na seção 1.3 do 
sesguinte livro:\cite{knuth1997}.


O Código 2 apresenta a implementação completa desta solução.

\begin{lstlisting}[language=C, caption={Troca sem variável auxiliar}, label={lst:q2}]
/* ====================================================
   Estrategia:
   - Usa apenas operacoes de soma e subtracao sobre
     as proprias variaveis a e b.
   - Passo 1: a = a + b  ->  a agora guarda a soma dos dois
   - Passo 2: b = a - b  ->  subtrai b da soma, restando
                              o valor original de a
   - Passo 3: a = a - b  ->  subtrai o novo b (antigo a)
                              da soma, restando o valor
                              original de b
   Exemplo com a=5, b=3:
     a = 5+3 = 8
     b = 8-3 = 5   (valor original de a)
     a = 8-5 = 3   (valor original de b)
   ==================================================== */

#include <stdio.h>

int main() {
    int a, b;

    printf("Digite dois inteiros: ");
    scanf("%d %d", &a, &b);

    printf("Antes: a = %d, b = %d\n", a, b);

    a = a + b;
    b = a - b;
    a = a - b;

    printf("Depois: a = %d, b = %d\n", a, b);

    return 0;
}
\end{lstlisting}

\section*{Conclusão}


O presente trabalho abordou o problema clássico de troca de valores entre duas 
variáveis inteiras na linguagem C, apresentando duas soluções distintas que, 
embora conduzam ao mesmo resultado, diferem em estratégia algoritmica e nos 
conceitos que mobilizam.

       
A primeira solução, baseada em uma variável auxiliar, destaca-se pela clareza 
e segurança. Sua lógica é intuitiva e de fácil compreensão, o que a torna a 
abordagem recomendada na maioria dos cenários práticos. A introdução de uma 
terceira variável como espaço temporário de armazenamento ilustra de forma 
didática o funcionamento do comando de atribuição e o princípio de que uma 
atribuição sobrescreve o valor anterior de uma variável.


A segunda solução, que dispensa memória adicional ao utilizar exclusivamente 
operações de soma e subtração, demonstra como propriedades aritméticas podem 
ser exploradas para resolver problemas de manipulação de dados. Embora elegante, 
essa abordagem apresenta a limitação de possível "overflow" quando os valores 
envolvidos são muito grandes, o que deve ser considerado em aplicações reais.


Ambas as implementações reforçam conceitos fundamentais trabalhados na 
disciplina de Algoritmos I, como declaração de variáveis, operador de 
atribuição, entrada e saída de dados com "scanf" e "printf", e construção de 
expressões aritméticas. A análise comparativa das duas abordagens contribui 
para o desenvolvimento do pensamento computacional, evidenciando que um mesmo 
problema pode admitir múltiplas soluções com diferentes compromissos entre 
simplicidade, eficiência e robustez.

\bibliographystyle{abntex2-alf}  % ou plain, ieeetr, etc.
\bibliography{references}

\end{document}
